% !TEX TS-program = pdflatex
% !TEX encoding = UTF-8 Unicode

% This file is a template using the "beamer" package to create slides for a talk or presentation
% - Giving a talk on some subject.
% - The talk is between 15min and 45min long.
% - Style is ornate.

% MODIFIED by Jonathan Kew, 2008-07-06
% The header comments and encoding in this file were modified for inclusion with TeXworks.
% The content is otherwise unchanged from the original distributed with the beamer package.

\documentclass{beamer}
\setbeamercovered{invisible}



% Copyright 2004 by Till Tantau <tantau@users.sourceforge.net>.
%
% In principle, this file can be redistributed and/or modified under
% the terms of the GNU Public License, version 2.
%
% However, this file is supposed to be a template to be modified
% for your own needs. For this reason, if you use this file as a
% template and not specifically distribute it as part of a another
% package/program, I grant the extra permission to freely copy and
% modify this file as you see fit and even to delete this copyright
% notice. 


\mode<presentation>
{
  \usetheme{Warsaw}
  % or ...

  % or whatever (possibly just delete it)
}


\usepackage[english]{babel}
% or whatever

\usepackage[utf8]{inputenc}
% or whatever

\usepackage{times}
\usepackage[T1]{fontenc}
% Or whatever. Note that the encoding and the font should match. If T1
% does not look nice, try deleting the line with the fontenc.

%%% MATH RELATED

\input{macros.tex}

\title{MATH 350-2 Advanced Calculus} 
\subtitle
{} % (optional)

\author[W.R. Casper] % (optional, use only with lots of authors)
{W.R. Casper}
% - Use the \inst{?} command only if the authors have different
%   affiliation.

\institute[California State University Fullerton] % (optional, but mostly needed)
{
  Department of Mathematics\\
  California State University Fullerton}
% - Use the \inst command only if there are several affiliations.
% - Keep it simple, no one is interested in your street address.

\subject{Talks}
% This is only inserted into the PDF information catalog. Can be left
% out. 



% If you have a file called "university-logo-filename.xxx", where xxx
% is a graphic format that can be processed by latex or pdflatex,
% resp., then you can add a logo as follows:

% \pgfdeclareimage[height=0.5cm]{university-logo}{university-logo-filename}
% \logo{\pgfuseimage{university-logo}}



% Delete this, if you do not want the table of contents to pop up at
% the beginning of each subsection:
\AtBeginSubsection[]
{
  \begin{frame}<beamer>{Outline}
    \tableofcontents[currentsection,currentsubsection]
  \end{frame}
}


% If you wish to uncover everything in a step-wise fashion, uncomment
% the following command: 

%\beamerdefaultoverlayspecification{<+->}


\begin{document}

\begin{frame}
  \titlepage
\end{frame}

\begin{frame}{Outline}
  \tableofcontents
  % You might wish to add the option [pausesections]
\end{frame}

% Since this a solution template for a generic talk, very little can
% be said about how it should be structured. However, the talk length
% of between 15min and 45min and the theme suggest that you stick to
% the following rules:  

% - Exactly two or three sections (other than the summary).
% - At *most* three subsections per section.
% - Talk about 30s to 2min per frame. So there should be between about
%   15 and 30 frames, all told.

\section{Real Analysis Lecture 20}

\subsection{More on differentiation}

\begin{frame}{Local maxima and minima}
\begin{defn}
Let $(M,d)$ be a metric space and $f$ be a real-valued function on a subset $S\subseteq M$.
\pause
Then $f$ has an \textbf{absolute (or global) maximum at $c\in S$} if
$$f(x)\leq f(c)\quad\text{for all}\ x\in S.$$
\pause
It has \textbf{local maximum at $a\in S$} if there exists an $r>0$ such that
$$f(x)\leq f(c)\quad\text{for all}\ x\in B_M(c;r)\cap S.$$
\pause
Absolute and local minima are defined similarly.
\end{defn}
\end{frame}

\begin{frame}{Critical points}
\begin{defn}
Let $f$ be a real-valued function on an open interval $(a,b)$.
\pause
A point $c\in (a,b)$ is called a \textbf{critical point} of $f$ if either $f$ is not differentiable at $c$ or else $f'(c) = 0$.
\end{defn}
\pause
\begin{thm}
Let $f$ be a real-valued function on an open interval $(a,b)$ and suppose that $f$ has a local maximum or local minimum at $c\in (a,b)$.
Then $c$ is a critical point of $f$.
\end{thm}
\end{frame}

\begin{frame}{Critical points}
\begin{proof}
\pause
We will prove this in the case $c$ is a local maximum.  The proof in the case of a local minimum is similar.\\
\pause
If $f$ is not differentiable at $c$, then we're done!\\
\pause
Therefore assume that it is differentiable.\\
\pause
By definition, there exists $r>0$ such that
$$f(x)\leq f(c)\quad\text{for all}\ x\in (c-r,c+r)\cap (a,b).$$
\pause
To obtain a contradiction, suppose that $f'(c)\neq 0$.\\
\pause
Then by the Helper Theorem, there is a function $f^*(x)$ defined on $(a,b)$ and continuous at $c\in (a,b)$ with
$$f(x)-f(c) = (x-c)f^*(x).$$
\end{proof}
\end{frame}

\begin{frame}{Critical points}
\begin{proof}
Moreover, $f^*(c) = f'(c)\neq 0$.\\
\pause
Since $f$ is continuous, we may choose $\delta > 0$ such that
$$|x-c| < \delta\ \ \Rightarrow\ \ |f^*(x) - f^*(c)| < |f^*(c)|/2.$$
\pause
Moreover, we may assume that $\delta < r$.\\
\pause
It follows that $f^*(x) > |f^*(c)|/2$ for all $x\in (a,b)$ with $|x-c|<\delta$.
\pause
However,
$$f(c\pm \delta/2) = f(c) \pm \frac{1}{2}\delta f^*(c+\delta/2),$$
one of these will be larger than $f(c)$.  \pause Contradiction!
\end{proof}
\end{frame}

\begin{frame}{Rolle's Theorem}
\begin{thm}{Rolle's Theorem}
Suppose that $f(x)$ is a real-valued function on $[a,b]$ which is
\begin{itemize}
\pause
\item continuous on $[a,b]$;
\pause
\item differentiable at each point in $(a,b)$;
\pause
\item and which satisfies $f(a) = f(b)$.
\end{itemize}
\pause
Then there exists $c\in (a,b)$ with $f'(c) = 0$.
\end{thm}
\end{frame}

\begin{frame}{Rolle's Theorem}
\begin{proof}
Suppose otherwise.\\
\pause
By the Extreme Value Theorem, $f(x)$ has an absolute maximum on $[a,b]$, as well as an absolute minimum on $[a,b]$.\\
\pause
If both of these occur at the endpoints, then the absolute max is the absolute min, forcing $f$ to be constant.\\
\pause
Then $f'(c)=0$ for every $c\in (a,b)$.\\
\pause
Otherwise at least one of the points happens inside, so $f(x)$ has a local max or local min at a value $c$ inside $(a,b)$.\\
\pause
The point $c$ is a critical point, and since $f$ is differentiable at $c$, we must have $f'(c)=0$.
\end{proof}
\end{frame}

\begin{frame}{Mean Value Theorem}
\begin{thm}{Rolle's Theorem}
Suppose that $f(x)$ is a real-valued function on $[a,b]$ which is
\begin{itemize}
\pause
\item continuous on $[a,b]$;
\pause
\item and differentiable at each point in $(a,b)$.
\end{itemize}
\pause
Then there exists $c\in (a,b)$ with
$$f'(c) = \frac{f(b)-f(a)}{b-a}.$$
\end{thm}
\end{frame}

\begin{frame}{Mean Value Theorem}
\begin{proof}
Consider the function
$$g(x) = f(x)-\left(\frac{f(b)-f(a)}{b-a}\right)(x-a) + f(a).$$
\pause
Then $g(x)$ is
\begin{itemize}
\pause
\item continuous on $[a,b]$;
\pause
\item differentiable at each point in $(a,b)$;
\pause
\item and satisfies $g(a) = g(b)$.
\end{itemize}
\pause
By Rolle's Theorem, there exists $c\in (a,b)$ with
$$0 = g'(c) = f'(c) -\frac{f(b)-f(a)}{b-a}.$$
\end{proof}
\end{frame}

\subsection{Riemann-Stieltjes integrals}

\begin{frame}{Integrals}
Back in Calculus, we learned about \textbf{integrals}.\\
\pause
It turns out that there are many types of integrals:
\begin{itemize}
\pause
\item the Riemann integral (our familiar friend)
\pause
\item the Riemann-Stieltjes integral
\pause
\item the Lebesgue integral ...
\end{itemize}
\pause
Most of these generalize the Riemann integral in some sense.
\end{frame}

\begin{frame}{Partitions}
\begin{defn}
A \textbf{partition} $P$ of $[a,b]$ Is a set of points
$$\{x_0,x_1,\dots,x_n\}$$
with
$$a=x_0 < x_1 < x_2 < \dots < x_n =b.$$
\pause
A set of \textbf{sample points} for $P$ is a set $\{t_1,\dots, t_n\}$ with $t_k\in [x_{k-1},x_k]$ for all $1\leq k\leq n$.
\pause
A partition $P'$ is a \textbf{refinement} of $P$ if $P\subseteq P'$.
\end{defn}
\end{frame}

\begin{frame}{Riemann-Stieltjes sums}
Notation: given a real-valued function $\alpha$ on $[a,b]$\\
$$\Delta \alpha_k(P) = \alpha(x_k)-\alpha(x_{k-1}),\quad k=1,2,\dots,n.$$
\pause
\begin{defn}
\pause
Let $f$ and $\alpha$ be real-valued functions on $[a,b]$, and let $P = \{x_0,x_1,\dots,x_n\}$ be a partition of $[a,b]$.
\pause
A sum of the form
$$S(P,f,\alpha,\{t_k\}) = \sum_{k=1}^n f(t_k) \Delta\alpha_k(P),$$
for some set of sample points $\{t_k\}$ of $P$.
\end{defn}
\end{frame}

\begin{frame}{Riemann-Stieltjes integral}
\begin{defn}
\pause
Let $f$ and $\alpha$ be real-valued functions on $[a,b]$.
\pause
We say $f$ is \textbf{Riemann integrable} with respect to $\alpha$ on $[a,b]$ if there exists $A\in\mathbb{R}$ with the following property:
\pause 
for all $\epsilon > 0$, there exists a partition $P_\epsilon$ of $[a,b]$ such that for all refinements $P$ of $P_\epsilon$
$$|S(P,f,\alpha,\{t_k\}) - A| < \epsilon,$$
for all sets of sample points $\{t_k\}$ of $P$.
\end{defn}
\pause
Notation: in this case $A$ is denote by $\int_a^b f d\alpha$ or $\int_a^b f(x) d\alpha(x)$
\end{frame}

\begin{frame}{Linearity of integration}
\begin{thm}
Suppose that $f$ and $g$ are Riemann integrable with respect to $\alpha$ on $[a,b]$.\\
\pause
Then for any $c_1,c_2\in\mathbb{R}$, $c_1f + c_2g$ is Riemann integrable with respect to $\alpha$ and
$$\int_a^b (c_1f + c_2g) d\alpha = c_1\int_a^b f d\alpha + c_2\int_a^b g d\alpha.$$
\end{thm}

\end{frame}
\begin{frame}{Linearity of integration}
\begin{proof}
\pause
Let $\epsilon > 0$ and WLOG assume $c_1,c_2\neq 0$.\\
\pause
Choose a partition $P_\epsilon'$ such that
$$|S(P,f,\alpha,\{t_k\}) - \int_a^bf d\alpha| < \epsilon/2|c_1|,$$
for all refinements $P$ of $P_\epsilon'$ and sets of sample points $\{t_k\}$ of $P$.\\
\pause
Also, choose a partition $P_\epsilon''$ such that
$$|S(P,g,\alpha,\{t_k\}) - \int_a^bg d\alpha| < \epsilon/2|c_2|,$$
for all refinements $P$ of $P_\epsilon''$ and sets of sample points $\{t_k\}$ of $P$.\\
\pause
Let $P_\epsilon = P_\epsilon'\cup P_\epsilon''$. 
\pause
Then any refinement $P$ of $P_\epsilon$, is a refinement of $P_\epsilon'$ and $P_\epsilon''$.
\end{proof}
\end{frame}
\begin{frame}{Linearity of integration}
\begin{proof}
Therefore if $P$ is a refinement of $P_\epsilon$, and $\{t_k\}$ is a set of sample points of $P$, then by the triangle inequality
\pause
\begin{align*}
|S(P,c_1f+c_2g,\alpha,\{t_k\}) - (c_1\int_a^b fd\alpha+c_2\int_a^bgd\alpha)|\\
 = |c_1S(P,f,\alpha,\{t_k\}) + c_2S(P,g,\alpha,\{t_k\})- (c_1\int_a^b fd\alpha+c_2\int_a^bgd\alpha)|\\
 = |c_1S(P,f,\alpha,\{t_k\}) - c_1\int_a^b fd\alpha + c_2S(P,g,\alpha,\{t_k\})- c_2\int_a^bgd\alpha|\\
 \leq |c_1|\cdot |S(P,f,\alpha,\{t_k\}) - \int_a^b fd\alpha| + |c_2|\cdot |S(P,g,\alpha,\{t_k\})- \int_a^bgd\alpha|\\
 < \epsilon/2 + \epsilon/2 = \epsilon.
\end{align*}
\pause
Since $\epsilon>0$ was arbitrary, this proves the theorem.
\end{proof}
\end{frame}

\begin{frame}{Linearity of integration}
\begin{thm}
Suppose that $f$ is Riemann integrable with respect to both $\alpha$ and $\beta$ on $[a,b]$.\\
\pause
Then for any $c_1,c_2\in\mathbb{R}$, $f$ is Riemann integrable with respect to $c_1\alpha + c_2\beta$ and
$$\int_a^b f d(c_1\alpha + c_2\beta) = c_1\int_a^b f d\alpha + c_2\int_a^b f d\beta.$$
\end{thm}
\end{frame}

\begin{frame}{Challenge!}
\begin{prob}
Prove the last theorem!
\end{prob}
\end{frame}

\begin{frame}{Linearity of integration}
\begin{thm}
Suppose that $f$ is a real-valued function on $[a,b]$ and that $c\in (a,b)$.
\pause
If $\int_a^c fd\alpha$ and $\int_b^c fd\alpha$ exist, then $\int_a^cfd\alpha$ exists and
$$\int_a^c fd\alpha + \int_c^b fd\alpha = \int_a^b fd\alpha.$$
\pause
Conversely, if $\int_a^c fd\alpha$ exists then $\int_a^c fd\alpha$ and $\int_b^c fd\alpha$ exist and
$$\int_a^c fd\alpha + \int_c^b fd\alpha = \int_a^b fd\alpha.$$
\end{thm}
\end{frame}


\end{document}


